\documentclass[11pt]{article}
\usepackage[a4paper,margin=2cm]{geometry}
\usepackage[spanish]{babel}
\usepackage[utf8]{inputenc}
\usepackage[T1]{fontenc}
\usepackage{amsmath,amssymb}
\usepackage{booktabs}
\pagestyle{empty}
\setlength{\parindent}{0pt}
\setlength{\parskip}{4pt}

\begin{document}

\begin{center}
{\large\bfseries Física Computacional (106018C)} \\[2pt]
{\bfseries Taller 3 --- Ajuste lineal y péndulo simple} \\[2pt]
Semana 3 \quad --- \quad Lenguajes: Fortran/Python \\[2pt]
Cristian Andrés Cárdenas Muñoz \quad --- \quad 02/10/2026
\end{center}

\vspace{4pt}
\hrule
\vspace{6pt}

\textbf{1. Contexto físico.}
Para oscilaciones pequeñas, el período de un péndulo simple satisface $T=2\pi\sqrt{L/g}$. Al medir el período para distintas longitudes y linealizar la relación, es posible estimar experimentalmente la aceleración de la gravedad.

\textbf{2. Método.}
Se limpiaron los datos experimentales antes del ajuste: se corrigieron registros justificables, se restringió el análisis a $|\theta|\leq 10^{\circ}$ y se excluyeron mediciones con problemas experimentales documentados. En Fortran se realizó un ajuste lineal por mínimos cuadrados usando
\[
T^2=aL+b, \qquad a=\frac{N\sum xy-(\sum x)(\sum y)}{N\sum x^2-(\sum x)^2}, \qquad g=\frac{4\pi^2}{a}.
\]

\textbf{3. Resultado.}
El programa procesó $N=89$ mediciones y obtuvo $a=4.04873469\,\mathrm{s^2/m}$, $b=-0.00288786\,\mathrm{s^2}$, $R^2=0.99974782$ y $g=9.75080380\,\mathrm{m/s^2}$.

\begin{center}
\begin{tabular}{lccc}
\toprule
Caso & Valor obtenido & Valor de referencia & Error relativo \\
\midrule
$g$ & $9.7508\,\mathrm{m/s^2}$ & $9.81\,\mathrm{m/s^2}$ & $0.603\%$ \\
\bottomrule
\end{tabular}
\end{center}

El valor de $R^2$ cercano a 1 indica una relación lineal muy fuerte entre $T^2$ y $L$. El intercepto pequeño y negativo sugiere una desviación sistemática reducida respecto al modelo ideal $b=0$, mientras que los residuos permanecen alrededor de cero sin una desviación dominante evidente.

\textbf{Verificación (PASS/FAIL):} 11 de 11 pruebas pasaron.

\textbf{4. Obstáculo o decisión de implementación.}
Fue necesario distinguir entre errores corregibles y mediciones que debían excluirse. Por ejemplo, algunos tiempos correspondían al período de una oscilación y no al tiempo total; se corrigieron usando el número de oscilaciones registrado. Además, el programa se diseñó para leer un número desconocido de filas mediante \texttt{iostat}, evitando asumir que siempre existirían 100 registros.

\textbf{5. Conclusión.}
El ajuste lineal permitió recuperar un valor de $g$ con un error relativo menor al 1\%. La práctica mostró que la calidad del resultado depende tanto del algoritmo de mínimos cuadrados como de una limpieza de datos justificable y del respeto al dominio de validez del modelo de ángulo pequeño.

\textbf{6. Uso de herramientas de IA.}
Se utilizó ChatGPT como apoyo para interpretar las instrucciones, revisar la lógica del código, documentar decisiones y redactar el reporte. Los resultados se verificaron compilando el programa y mediante la celda automática PASS/FAIL del notebook.

\vspace{4pt}
\hrule
\vspace{4pt}
{\small Código fuente: ver archivo adjunto \texttt{ajuste\_pendulo.f90}.}

\end{document}
